\section{An initial simplex from a flag}\label{sec:flag}

We now apply the leading-term construction to the section ring of $L$.
A complete-intersection flag gives a particularly simple bound on the
exponents: a simplex whose volume equals the leading coefficient of the
Hilbert polynomial.  We give the complete argument for the flag
construction of \cite[Section~3]{OpenAI2026Higher}.

For positive real numbers $a_1,\ldots,a_d$, write
\[
 \Delta(a_1,\ldots,a_d)
 =\left\{x\in\R_{\geq0}^d:
       \sum_{i=1}^d\frac{x_i}{a_i}\leq1\right\}.
\]
Its volume is $\prod_i a_i/d!$.  We use additive notation $NL$ for
$L^{\otimes N}$.

\begin{proposition}\label{prop:initial-simplex}
Let $X$ be a smooth integral projective $n$-dimensional variety over an
algebraically closed field, with $n\geq2$, and let $L$ be ample.  Put
$V=L^n$, and choose an integer $D\geq1$ such that $DL$ is very ample.
There is a graded family of finite sets $S_N\subset\Z_{\geq0}^n$ such
that
\begin{align}
 S_N&\subset NP_0,
 &P_0&=\Delta(1/D,\ldots,1/D,D^{n-1}V),
 \label{eq:initial-simplex}\\
 \#S_N&=h^0(X,NL)
       =\frac{V}{n!}N^n+o(N^n).
 \label{eq:initial-dimension}
\end{align}
For every $N,r\geq1$ and $m\geq0$, if $M(S_N)$ separates jets through
order $m$ at some $r$ distinct torus points, then $H^0(X,NL)$ separates
the same jets at some $r$ distinct points of $X$.
\end{proposition}

\begin{proof}
Bertini's theorem for hyperplane sections, valid in every characteristic
for a projective embedding over an algebraically closed field, gives
sections $\sigma_1,\ldots,\sigma_{n-1}\in H^0(X,DL)$ whose successive
zero loci form a smooth integral flag
\[
 X=X_0\supset X_1\supset\cdots\supset X_{n-1}=C_0.
\]
See \cite[Tags~0FD6 and~0G4G]{StacksProject} for the smoothness and
integrality assertions.
At each step the variety being cut has dimension at least two.  Choose
$x\in C_0(k)$ and a local frame $e$ of $L$ near $x$.  The functions
$z_i=\sigma_i/e^D$, $i<n$, form part of a regular system of parameters
at $x$; complete them by $z_n$.  These parameters identify
$\widehat{\mathcal O}_{X,x}$ with $k[[z_1,\ldots,z_n]]$.

Expand every section of $NL$ using the frame $e^N$, and let $S_N$ be
the set of lexicographically least exponents of the nonzero expansions.
The expansion map is injective: a section with zero germ is zero by
integrality, and the local ring injects into its completion.
Lemma~\ref{lem:least-terms} therefore gives
$\#S_N=h^0(X,NL)$ and the graded property.  The Hilbert polynomial of
the ample line bundle gives the last equality in
\eqref{eq:initial-dimension}.  A basis with distinct least exponents,
Lemma~\ref{lem:exposing-weights}(i), and
Lemma~\ref{lem:formal-transfer} with $h=n$ prove the jet-transfer
assertion.

It remains to prove the simplex bound.  Let
$\alpha=(\alpha_1,\ldots,\alpha_n)\in S_N$, coming from a nonzero
section $s$.  We successively divide by the equations of the flag and
restrict to the next member.  The details matter because the
lexicographic order first measures vanishing along an entire divisor,
not just at the chosen point.

Suppose we have reached a nonzero section on $X_{i-1}$ with least
exponent $(\alpha_i,\ldots,\alpha_n)$ in its formal expansion at $x$.
If $\alpha_i>0$, restriction of this section to $X_i$ has zero formal
expansion.  Injectivity into the completed local ring, followed by
integrality of $X_i$, shows that this restriction is identically zero.
The exact sequence for the effective Cartier divisor
$X_i\subset X_{i-1}$ therefore allows division by
$\sigma_i|_{X_{i-1}}$.  Repeating this argument $\alpha_i$ times, and
then restricting to $X_i$, gives a nonzero section whose least exponent
is $(\alpha_{i+1},\ldots,\alpha_n)$.  Nonvanishing follows from the
nonzero coefficient of the original least term.  Each division
subtracts $DL$ from the line bundle.

After these steps, the resulting nonzero section on $C_0$ belongs to
\[
 \left(N-D\sum_{i<n}\alpha_i\right)L|_{C_0}
\]
and has order exactly $\alpha_n$ at $x$.  Since
$L\cdot C_0=D^{n-1}V$, comparison with the degree of its zero divisor
gives
\[
 \alpha_n\leq
 \left(N-D\sum_{i<n}\alpha_i\right)D^{n-1}V.
\]
Equivalently,
$D\sum_{i<n}\alpha_i+\alpha_n/(D^{n-1}V)\leq N$, which is
\eqref{eq:initial-simplex}.
\end{proof}

The intercepts of $P_0$ have product $V$.  The next section changes
them, two at a time, without changing this product, the numbers of
exponents, or backward jet transfer.
